%==============================================================================
% CAPÍTULO 9 — RADIOLOGIA ODONTOLÓGICA COM DEEP LEARNING
% Parte II: Teoria e Modelos de Deep Learning
%==============================================================================

\chapter{Radiologia Odontológica com Deep Learning}
\label{cap:radiologia-dl}

\nivelquatro

\begin{quote}
\itshape\small
``Deep learning-based caries detection in panoramic radiographs achieves pooled
sensitivity of 89.5\% and specificity of 91.3\%, demonstrating diagnostic
performance comparable to experienced clinicians.'' --- \citeonline{ODT-016},
Journal of Dental Research, 2024.
\end{quote}

A radiologia é a especialidade odontológica mais impactada pelo deep learning.
Da radiografia periapical ao CBCT, algoritmos de IA estão sendo integrados ao
fluxo de trabalho clínico para detecção de cáries, avaliação de qualidade óssea,
identificação de lesões periapicais e triagem de patologias. Este capítulo aborda
as três principais modalidades radiográficas --- panorâmica, bitewing/periapical
e CBCT --- e as arquiteturas de deep learning que as analisam.

\notaexplicativa{Este capítulo assume domínio de CNNs (Capítulo 7) e segmentação
(Capítulo 8). Para leitores sem formação em radiologia odontológica, recomenda-se
revisão dos princípios básicos de formação de imagem radiográfica antes de
prosseguir (kVp, mA, densidade radiográfica, geometria de projeção).}

\section{Objetivos de Aprendizagem}

Ao final deste capítulo, você será capaz de:

\subsection{Compreender as Modalidades Radiográficas Odontológicas}
Você aprenderá as características técnicas de panorâmicas, bitewing, periapicais e CBCT que impactam o projeto de modelos de deep learning.

\subsection{Detectar Cáries e Lesões com Deep Learning}
Exploraremos arquiteturas e datasets para detecção automatizada de cáries, lesões periapicais e avaliação de qualidade óssea.

\subsection{Aplicar Redução de Artefatos com Redes Neurais}
Você entenderá como a IA pode reconstruir imagens de baixa dose e remover artefatos metálicos em CBCT.

\subsection{Construir um Pipeline Completo de Diagnóstico Radiográfico}
Ao final, você implementará um sistema integrado de classificação e detecção para radiografias odontológicas.

%-----------------------------------------------------------------------------
\section{Modalidades Radiográficas Odontológicas}
\label{sec:modalidades}

\subsection{Radiografia Panorâmica}

A radiografia panorâmica (ortopantomografia) é o exame de triagem mais comum
na odontologia, capturando maxila, mandíbula, dentes, ATM e seios maxilares
em uma única imagem com dose efetiva de $\sim$15--25 $\mu$Sv.

\subsubsection{Características para Deep Learning}

\textbf{Características relevantes para DL:}
\begin{itemize}
  \item \textbf{Resolução típica:} $2400 \times 1200$ a $3000 \times 1500$ pixels
  \item \textbf{Distorção inerente:} Projeção curva gera alongamento vertical
        variável (fator de magnificação 1.2--1.3$\times$) e sobreposição de
        estruturas (coluna cervical, osso hioide). Modelos DL precisam aprender
        a ser invariantes a estas distorções.
  \item \textbf{Datasets públicos:} Tufts (1000+ imagens), UFBA-UESC (425 imagens
        com bounding boxes e polígonos), DNS (543 imagens de segmentação)
  \item \textbf{Tarefas DL:} Detecção de cáries, classificação de qualidade óssea,
        detecção de lesões periapicais, contagem de dentes, identificação de
        restaurações e implantes
\end{itemize}

\citacaoativa{doi-5051-jpis-2302880144}{10.5051/jpis.2302880144}{O deep learning aplicado a radiografias
panorâmicas para avaliação de qualidade óssea em 2.270 sítios edêntulos demonstrou
correlação significativa com CBCT ($r = 0.702$) e tato cirúrgico ($r = 0.658$)
utilizando ResNet-50 com fine-tuning (Lee et al., 2024).}

\subsection{Radiografia Bitewing e Periapical}

\begin{itemize}
  \item \textbf{Bitewing:} Projeção intraoral com o filme/sensor posicionado
        paralelo às coroas. Ideal para detecção de cáries interproximais e
        avaliação de crista óssea alveolar. Alta resolução ($\sim$20 lp/mm),
        baixa distorção.
  \item \textbf{Periapical:} Projeção intraoral que captura o dente completo
        (coroa até ápice). Técnica da bissetriz ou paralelismo. Essencial para
        endodontia (lesões periapicais, fraturas radiculares, reabsorções).
\end{itemize}

\textbf{Relevância para DL:} A geometria padronizada (paralelismo) das bitewings
as torna ideais para treinamento de CNNs, com menos variabilidade que panorâmicas.
\citeonline{ODT-017} (Cantu et al., 2023) demonstrou que ResNet-34 em bitewings
atingiu AUC 0.93 vs. 0.85 (média de 10 dentistas) na detecção de cáries proximais.

\subsection{CBCT (\textit{Cone Beam Computed Tomography})}

O CBCT revolucionou o diagnóstico 3D em odontologia com dose significativamente
menor que a TC médica convencional (30--100 $\mu$Sv vs. 2000+ $\mu$Sv).

\subsubsection{Características do CBCT para Deep Learning}

\begin{itemize}
  \item \textbf{Princípio:} Feixe cônico de raios-X gira em torno do paciente;
        detector \textit{flat-panel} captura projeções 2D; algoritmo de
        Feldkamp-Davis-Kress (FDK) reconstrói volume 3D.
  \item \textbf{Voxel:} Unidade volumétrica isotrópica (tipicamente 0.075--0.4 mm).
        Voxels menores (0.075 mm) permitem visualização de canais radiculares
        e trincas; voxels maiores (0.3--0.4 mm) reduzem dose e ruído.
  \item \textbf{Unidades Hounsfield (HU):} Diferentemente da TC médica, o CBCT
        não produz valores HU calibrados devido ao espalhamento, endurecimento
        de feixe e artefatos de reconstrução FDK. Isso limita a transferibilidade
        de modelos treinados em TC médica para CBCT.
  \item \textbf{Artefatos:} Metálicos (implantes, restaurações de amálgama),
        \textit{beam hardening} (faixas escuras entre objetos densos), ruído
        de baixa dose, \textit{cone-beam artifact} (distorção nas bordas do FOV).
        \citeonline{ODT-034} (ToothSegNet) simula degradação durante o treinamento
        para melhorar robustez.
  \item \textbf{Desafios DL:} Dimensionalidade 3D ($512^3$ voxels = 134M elementos),
        anotação custosa (segmentação manual de um CBCT leva 2--4 horas por
        especialista), variabilidade entre equipamentos.
\end{itemize}

\begin{table}[H]
  \centering
  \caption{Comparação das modalidades radiográficas para DL odontológico}
  \label{tab:modalidades-comp}
  \begin{tabularx}{\textwidth}{lccccX}
    \toprule
    \textbf{Característica} & \textbf{Panorâmica} & \textbf{Bitewing} & \textbf{Periapical} & \textbf{CBCT} & \textbf{Nota} \\
    \midrule
    Dimensionalidade & 2D & 2D & 2D & 3D & CBCT requer adaptação 3D \\
    Dose ($\mu$Sv) & 15--25 & 1--5 & 1--5 & 30--100 & DL não afeta dose \\
    Distorção & Moderada & Mínima & Variável & Isotrópica & Afeta augmentação \\
    Datasets públicos & 3--5 & 1--2 & 0--1 & 2--4 & Principal gargalo \\
    Custo anotação & Baixo & Baixo & Baixo & Muito alto & 2--4h/CBCT \\
    \bottomrule
  \end{tabularx}
\end{table}

%-----------------------------------------------------------------------------

% === FIGURAS ADICIONADAS (OdontoIA Dark) ===
% ======================================================================
% FIGURA: Matriz de Confusão para Classificação de Cárie (Capítulo 7/9)
% ======================================================================
\begin{figure}[H]
\centering
\begin{tikzpicture}[scale=0.9]
% Grid da matriz de confusão
\fill[codegreen!20] (0,0) rectangle (2,2);
\fill[red!20] (2,0) rectangle (4,2);
\fill[red!20] (0,-2) rectangle (2,0);
\fill[codegreen!20] (2,-2) rectangle (4,0);

% Linhas da grade
\draw[thick] (0,-2) -- (4,-2) -- (4,2) -- (0,2) -- cycle;
\draw (2,-2) -- (2,2);
\draw (0,0) -- (4,0);

% Labels
\node[align=center, font=\large\bfseries] at (1,1) {VP\\85};
\node[align=center, font=\large] at (3,1) {FP\\12};
\node[align=center, font=\large] at (1,-1) {FN\\8};
\node[align=center, font=\large\bfseries] at (3,-1) {VN\\95};

% Eixos
\node[font=\footnotesize] at (2,2.5) {\textbf{Real: Cárie}};
\node[font=\footnotesize] at (2,-2.8) {\textbf{Real: Sem Cárie}};
\node[font=\footnotesize, rotate=90] at (-1,1) {\textbf{Pred: Cárie}};
\node[font=\footnotesize, rotate=90] at (-1,-1) {\textbf{Pred: Sem Cárie}};

% Métricas
\node[font=\small, anchor=west] at (5,1.5) {Sensibilidade: $\frac{85}{93}=0.914$};
\node[font=\small, anchor=west] at (5,0.5) {Especificidade: $\frac{95}{107}=0.888$};
\node[font=\small, anchor=west] at (5,-0.5) {Acurácia: $\frac{180}{200}=0.900$};
\node[font=\small\bfseries, anchor=west] at (5,-1.5) {F1-Score: $0.895$};

\end{tikzpicture}
\caption{Matriz de confusão ilustrativa para um classificador de cárie dentária (Rohban et al., 2024). VP=Verdadeiro Positivo, FP=Falso Positivo, FN=Falso Negativo, VN=Verdadeiro Negativo. A diagonal principal (verde) representa acertos; a diagonal secundária (vermelho) representa erros.}
\label{fig:confusion-matrix}
\end{figure}

% ======================================================================
% FIGURA: Curva ROC comparativa (Capítulo 4/11)
% ======================================================================
\begin{figure}[H]
\centering
\begin{tikzpicture}[scale=0.85]
% Eixos
\draw[-latex, thick] (0,0) -- (7,0) node[right] {1 - Especificidade};
\draw[-latex, thick] (0,0) -- (0,7) node[above] {Sensibilidade};

% Linha de referência (random)
\draw[dashed, gray, thick] (0,0) -- (6,6) node[above left, font=\tiny] {AUC=0.50};

% Curva ROC do modelo bom (CNN)
\draw[thick, blueaccent] (0,0) .. controls (0.5,2) and (1,5.5) .. (2,6.2) -- (5,6.8) -- (6,7);
\node[font=\small, blueaccent] at (4.5,5.5) {CNN (AUC=0.93)};

% Curva ROC do modelo baseline (Random Forest)
\draw[thick, orangeaccent] (0,0) .. controls (0.5,1) and (2,3) .. (3,4.5) .. controls (4,5.2) and (5,5.5) .. (6,6.5);
\node[font=\small, orangeaccent] at (5.5,3.5) {RF (AUC=0.82)};

% Curva ROC do clínico
\draw[thick, codegreen] (0,0) .. controls (0.3,1.5) and (0.8,4) .. (1.5,5.5) .. controls (2,6) and (2.5,6.3) .. (6,6.9);
\node[font=\small, codegreen] at (2.5,5.0) {Clínico (AUC=0.88)};

% Grid
\draw[gray!30, thin] (1,0) -- (1,7);
\draw[gray!30, thin] (2,0) -- (2,7);
\draw[gray!30, thin] (3,0) -- (3,7);
\draw[gray!30, thin] (4,0) -- (4,7);
\draw[gray!30, thin] (5,0) -- (5,7);
\draw[gray!30, thin] (6,0) -- (6,7);

% Ticks
\foreach \x in {0,0.2,0.4,0.6,0.8,1.0} {
    \pgfmathsetmacro{\xp}{\x*6}
    \draw (\xp, -0.1) -- (\xp, 0.1);
    \node[font=\tiny, below] at (\xp, -0.2) {\x};
    \draw (-0.1, \xp) -- (0.1, \xp);
    \node[font=\tiny, left] at (-0.3, \xp) {\x};
}

\node[font=\footnotesize] at (3.5, -1) {Curvas ROC comparativas para diagnóstico de cárie (Rohban et al., 2024, meta-análise de 87 estudos)};
\end{tikzpicture}
\caption{Curvas ROC comparativas para detecção de cárie em radiografias panorâmicas. A CNN (azul, AUC=0.93) supera o clínico humano (verde, AUC=0.88) e o Random Forest baseline (laranja, AUC=0.82). Dados da meta-análise de Rohban et al. (2024).}
\label{fig:roc-curves}
\end{figure}


\section{Detecção de Cáries com Deep Learning}
\label{sec:caries-dl}

A cárie dentária é a doença crônica mais prevalente globalmente (afetando 2.4
bilhões de pessoas -- GBD 2017). A detecção precoce via radiografia é o padrão-ouro
para cáries interproximais, e o DL oferece um ``segundo leitor'' consistente.

\subsection{Arquiteturas Proeminentes}

\subsubsection{CariesNet~\cite{zhu2022cariesnet} (Zhu et al., 2022)}

Arquitetura focada em detecção multi-estágio de cáries em radiografias panorâmicas
com classificação de profundidade (esmalte, dentina, pulpar):

\begin{itemize}
  \item \textbf{Estágio 1:} Detecção de ROI (região dental) com Faster R-CNN
  \item \textbf{Estágio 2:} Classificação de profundidade da cárie com
        DenseNet-121 (3 classes: superficial, média, profunda)
  \item \textbf{Métrica:} AUC 0.92 para cárie em esmalte; 0.96 para cárie em dentina
\end{itemize}

\subsubsection{DCDNet~\cite{dayi2023dcdnet} (Dayı et al., 2023)}

Dual-Channel Deep Network para detecção simultânea de cáries proximais e oclusais:

\begin{itemize}
  \item \textbf{Canal 1:} CNN 2D processa a radiografia completa (contexto global)
  \item \textbf{Canal 2:} CNN 2D processa patches de alta resolução centrados
        em cada dente (detalhe local)
  \item \textbf{Fusão:} Concatenação tardia dos embeddings + classificador FC
  \item \textbf{Performance:} Acurácia 92.1\%, sensibilidade 89.7\%, especificidade 93.5\%
\end{itemize}

\subsection{Detecção Específica por Tipo Radiográfico}

\begin{enumerate}
  \item \textbf{Panorâmica:} \citeonline{ODT-016} (Lee et al., 2024) --
        metanálise de 22 estudos: sensibilidade 89.5\%, especificidade 91.3\%.
        U-Net obteve melhor Dice (0.87) para segmentação; YOLO melhor relação
        velocidade-acurácia.

  \item \textbf{Bitewing:} \citeonline{ODT-017} (Cantu et al., 2023) --
        ResNet-34 superou 10 dentistas (AUC 0.93 vs 0.85). CNN foi mais sensível
        para cáries iniciais (D1-D2), o estágio mais desafiador para clínicos.

  \item \textbf{Periapical (cárie radicular):} \citeonline{kim2023caries} (Kim et al.,
        2023) -- DenseNet-121, acurácia 91.7\%, sensibilidade 89.5\%,
        especificidade 93.2\% para cáries radiculares. Primeiro estudo dedicado
        a este tipo específico.

  \item \textbf{Comparação de arquiteturas:} \citeonline{ODT-022} (Carvalho et al.,
        2024) -- 7 arquiteturas comparadas em 8.500 bitewings. EfficientNet-B3
        obteve melhor equilíbrio (AUC 0.94, F1 0.91); MobileNetV2 foi 5$\times$
        mais rápido com apenas 2\% de perda de acurácia.
\end{enumerate}

\citacaoativa{rohban2024}{10.1016/j.jdent.2024.105398}{A meta-análise de 87 estudos
(Rohban et al., 2024) confirma que CNNs para detecção de cáries atingem
sensibilidade agrupada de 0.87 e especificidade de 0.91 em bitewing.
Transfer learning e aumento de dados são fatores críticos de sucesso;
arquiteturas treinadas do zero consistentemente subperformam.}

%-----------------------------------------------------------------------------
\section{Classificação de Qualidade Óssea}
\label{sec:qualidade-ossea}

A classificação da qualidade/tipo ósseo é crucial para planejamento de implantes
(osso tipo IV tem taxa de falha 35\% maior que tipo I-II). Tradicionalmente
avaliada pelo tato cirúrgico (Lekholm \& Zarb, 1985) ou CBCT (análise de
densidade), a IA oferece uma alternativa objetiva e não-invasiva.

\subsection{Estudos Seminais}

\subsubsection{Lee et al. (2024)}

\textbf{Lee et al. (2024):} \citeonline{DL-ODT-002} -- ResNet-50 para classificação
óssea em panorâmicas:
\begin{itemize}
  \item 2.270 sítios edêntulos; classificação L\&Z (D1--D4)
  \item Correlação com CBCT: $r = 0.702$; com tato cirúrgico: $r = 0.658$
  \item AUC por classe: D1=0.78, D2=0.81, D3=0.74, D4=0.72
  \item Conclusão: panorâmica + DL oferece screening razoável; CBCT permanece
        necessário para planejamento definitivo
\end{itemize}

\subsubsection{Pornvoranant et al. (2024)}

\textbf{Pornvoranant et al. (2024):} \citeonline{DL-ODT-003} -- Comparação direta
IA vs. humanos em CBCT:
\begin{itemize}
  \item 5 modelos DL (Inception-ResNet-v2 líder) vs. 23 residentes + 2 implantologistas
  \item Acurácia DL: 86\% vs. pooled humano: $\sim$70\%
  \item DL superou residentes em todos os níveis de experiência; foi comparável
        aos implantologistas seniores
  \item Principal limitação: 1.100 slices de CBCT (dataset pequeno para os padrões DL)
\end{itemize}

\subsection{Pipeline de Classificação Óssea}

\begin{enumerate}
  \item \textbf{Localização do sítio:} Detecção da região edêntula na panorâmica ou
        seleção do corte transversal no CBCT
  \item \textbf{Extração de ROI:} Patch $64\times64$ ou $128\times128$ centrado
        no sítio de interesse
  \item \textbf{Classificação:} CNN (ResNet-50, EfficientNet) $\to$ 4 classes (D1-D4)
  \item \textbf{Interpretação:} Grad-CAM para validação clínica da região de
        interesse do modelo
\end{enumerate}

%-----------------------------------------------------------------------------
\section{Lesões Periapicais e Patologias}
\label{sec:lesoes-periapicais}

\subsection{Detecção de Lesões Periapicais}

\subsubsection{ConvNeXt para Lesões Periapicais}

\textbf{ConvNeXt para lesões periapicais~\cite{liu2024convnext} (Liu et al., 2024):}
\begin{itemize}
  \item Arquitetura ConvNeXt (modernizada da ResNet com design Swin-inspired)
  \item Dataset: 4.200 panorâmicas com anotações de lesões periapicais
  \item AUC 0.96, sensibilidade 92\%, especificidade 94\%
  \item Modelo particularmente eficaz para lesões pequenas ($<$ 3 mm) -- onde
        dentistas generalistas têm menor acurácia
\end{itemize}

\subsubsection{Grad-CAM para Explicabilidade}

\textbf{Grad-CAM para explicabilidade (Kim et al., 2024):}
\citeonline{DL-ODT-021} aplicou EfficientNet-B4 com Grad-CAM para detecção de
patologia periapical em 2.100 panorâmicas:
\begin{itemize}
  \item AUC 0.94; Grad-CAM validado por 3 radiologistas ($\kappa = 0.81$)
  \item Modelo identificou 91\% das lesões sutis (radiolucências periapicais
        $<$ 2 mm) -- performance superior à inspeção visual sem auxílio
  \item Importante: Grad-CAM confirmou que o modelo focava no ápice radicular
        e lâmina dura, não em artefatos ou sobreposições
\end{itemize}

\subsection{Detecção de Restaurações}

\subsubsection{YOLOv5 para Detecção de Restaurações}

\textbf{YOLOv5 para restaurações (Bayrakdar et al., 2024):} \citeonline{DL-ODT-005}
\begin{itemize}
  \item 1.200 panorâmicas com 4.500 anotações de restaurações
  \item YOLOv5x: mAP@0.5 = 0.94, precisão 0.93, recall 0.91
  \item Detectou restaurações metálicas (amálgama), cerâmicas, resinosas e coroas
  \item Aplicação: triagem automatizada em larga escala (saúde pública) e
        auditoria de qualidade de tratamentos prévios
\end{itemize}

%-----------------------------------------------------------------------------
\section{Redução de Artefatos e Reconstrução com IA}
\label{sec:artefatos-dl}

\subsection{O Problema dos Artefatos em CBCT}

Artefatos são a principal limitação do CBCT para diagnóstico fino. Diferentes
tipos de artefato exigem diferentes estratégias de mitigação por DL:

\begin{enumerate}
  \item \textbf{Artefatos metálicos (\textit{metal artifacts}):} Causados por
        implantes, coroas metálicas, amálgama. Produzem faixas hiperdensas
        (brancas) e hipodensas (escuras) que obscurecem estruturas adjacentes.
        \textbf{Abordagem DL:} Conditional GAN (pix2pix) para tradução
        imagem-com-artefato $\to$ imagem-sem-artefato. \citeonline{DL-ODT-015}
        (Liu et al., 2024) reportou melhoria de PSNR de 8.2 dB e redução de
        RMSE de 62\% usando cGAN com generator U-Net e discriminator PatchGAN.

  \item \textbf{\textit{Beam hardening}:} Endurecimento do feixe de raios-X ao
        atravessar estruturas densas, causando faixas escuras (``cupping'')
        entre objetos densos. \textbf{Abordagem DL:} Pré-processamento com CNN
        para correção do espectro de energia antes da reconstrução FDK.

  \item \textbf{Ruído de baixa dose:} Protocolos de baixa dose (ALARA --
        \textit{As Low As Reasonably Achievable}) reduzem radiação mas aumentam
        ruído estatístico. \textbf{Abordagem DL:} Denoising com autoencoders
        convolucionais ou ESRGAN (\citeonline{DL-ODT-016}) -- super-resolução
        e denoising simultâneos com RRDB (\textit{Residual-in-Residual Dense
        Blocks}).

  \item \textbf{\textit{Cone-beam artifact}:} Distorção nas bordas do FOV
        (campo de visão) devido à geometria do feixe cônico. \textbf{Abordagem
        DL:} Extrapolação de FOV -- CNN aprende a ``completar'' estruturas
        parcialmente capturadas nas bordas.
\end{enumerate}

\subsection{Pipeline de Redução de Artefatos com cGAN}

\begin{lstlisting}[language=Python, caption={Redução de artefatos metálicos em
CBCT com Conditional GAN}, label={lst:cgan-cbct}]
import tensorflow as tf
from tensorflow.keras import layers, Model

# Generator: U-Net (traduz CBCT com artefato -> CBCT limpo)
def build_artifact_removal_generator(input_shape=(256, 256, 1)):
    inputs = layers.Input(input_shape)

    # Encoder
    e1 = layers.Conv2D(64, 4, strides=2, padding='same')(inputs)
    e1 = layers.LeakyReLU(0.2)(e1)
    e2 = layers.Conv2D(128, 4, strides=2, padding='same')(e1)
    e2 = layers.BatchNormalization()(e2)
    e2 = layers.LeakyReLU(0.2)(e2)
    e3 = layers.Conv2D(256, 4, strides=2, padding='same')(e2)
    e3 = layers.BatchNormalization()(e3)
    e3 = layers.LeakyReLU(0.2)(e3)

    # Bottleneck
    b = layers.Conv2D(512, 4, strides=2, padding='same')(e3)
    b = layers.LeakyReLU(0.2)(b)

    # Decoder com skip connections
    d1 = layers.Conv2DTranspose(256, 4, strides=2, padding='same')(b)
    d1 = layers.BatchNormalization()(d1)
    d1 = layers.Dropout(0.5)(d1)
    d1 = layers.Concatenate()([d1, e3])

    d2 = layers.Conv2DTranspose(128, 4, strides=2, padding='same')(d1)
    d2 = layers.BatchNormalization()(d2)
    d2 = layers.Concatenate()([d2, e2])

    d3 = layers.Conv2DTranspose(64, 4, strides=2, padding='same')(d2)
    d3 = layers.BatchNormalization()(d3)
    d3 = layers.Concatenate()([d3, e1])

    outputs = layers.Conv2DTranspose(1, 4, strides=2, padding='same',
                                     activation='tanh')(d3)

    return Model(inputs, outputs)

# Discriminator: PatchGAN (classifica patches como reais/falsos)
def build_patchgan_discriminator(input_shape=(256, 256, 2)):
    """Entrada: concatenação de [imagem_com_artefato, imagem_limpa_ou_gerada]"""
    inputs = layers.Input(input_shape)

    x = layers.Conv2D(64, 4, strides=2, padding='same')(inputs)
    x = layers.LeakyReLU(0.2)(x)
    x = layers.Conv2D(128, 4, strides=2, padding='same')(x)
    x = layers.BatchNormalization()(x)
    x = layers.LeakyReLU(0.2)(x)
    x = layers.Conv2D(256, 4, strides=2, padding='same')(x)
    x = layers.BatchNormalization()(x)
    x = layers.LeakyReLU(0.2)(x)
    x = layers.Conv2D(512, 4, strides=1, padding='same')(x)
    x = layers.BatchNormalization()(x)
    x = layers.LeakyReLU(0.2)(x)
    outputs = layers.Conv2D(1, 4, strides=1, padding='same')(x)

    return Model(inputs, outputs)

# Loss functions
cross_entropy = tf.keras.losses.BinaryCrossentropy(from_logits=True)
def generator_loss(disc_generated_output, gen_output, target, lambda_l1=100):
    gan_loss = cross_entropy(tf.ones_like(disc_generated_output),
                             disc_generated_output)
    l1_loss = tf.reduce_mean(tf.abs(target - gen_output))
    return gan_loss + lambda_l1 * l1_loss, gan_loss, l1_loss

def discriminator_loss(disc_real_output, disc_generated_output):
    real_loss = cross_entropy(tf.ones_like(disc_real_output), disc_real_output)
    generated_loss = cross_entropy(tf.zeros_like(disc_generated_output),
                                   disc_generated_output)
    return real_loss + generated_loss

# Treinamento (loop manual para GAN)
@tf.function
def train_step(input_image, target_image, generator, discriminator,
               gen_optimizer, disc_optimizer):
    with tf.GradientTape() as gen_tape, tf.GradientTape() as disc_tape:
        gen_output = generator(input_image, training=True)

        disc_real_output = discriminator(
            tf.concat([input_image, target_image], axis=-1), training=True)
        disc_generated_output = discriminator(
            tf.concat([input_image, gen_output], axis=-1), training=True)

        gen_total_loss, gen_gan_loss, gen_l1_loss = generator_loss(
            disc_generated_output, gen_output, target_image)
        disc_loss = discriminator_loss(disc_real_output, disc_generated_output)

    gen_gradients = gen_tape.gradient(gen_total_loss, generator.trainable_variables)
    disc_gradients = disc_tape.gradient(disc_loss, discriminator.trainable_variables)

    gen_optimizer.apply_gradients(zip(gen_gradients, generator.trainable_variables))
    disc_optimizer.apply_gradients(
        zip(disc_gradients, discriminator.trainable_variables))

    return gen_total_loss, disc_loss

print("Gerador e Discriminador criados. Pronto para treinamento GAN.")
print("Pipeline: CBCT com artefato -> Generator (U-Net) -> CBCT limpo")
\end{lstlisting}

\subsection{Geração de Imagens Sintéticas com StyleGAN3}

Além de remover artefatos, GANs podem \textbf{criar} imagens sintéticas para
aumentar datasets odontológicos escassos. \citeonline{DL-ODT-014} (Mehrfar et
al., 2024) utilizaram StyleGAN3 para gerar 5.000 radiografias panorâmicas
sintéticas a partir de 800 reais:

\begin{itemize}
  \item \textbf{Avaliação de qualidade:} Dentistas não distinguiram imagens
        sintéticas de reais em teste de Turing ($p > 0.05$)
  \item \textbf{Impacto no downstream:} CNN detectora de cáries treinada com
        dados reais + sintéticos melhorou F1-score de 0.72 para 0.78 (+8.3\%)
  \item \textbf{Vantagem ética:} Geração de dados sintéticos não requer
        consentimento de paciente e pode balancear datasets (ex.: gerar mais
        exemplos de lesões raras como eritroplasia)
  \item \textbf{Cuidado:} \textit{Mode collapse} -- a GAN pode gerar apenas
        variações limitadas. Monitorar diversidade com FID (\textit{Fréchet
        Inception Distance}); valores $< 50$ indicam boa diversidade
\end{itemize}

\citacaoativa{doi-1016-j-compmedimag-2024-102347}{10.1093/dmfr/twae044}{A geração de radiografias
panorâmicas sintéticas com StyleGAN3 (5.000 imagens a partir de 800 reais)
melhorou o F1-score da CNN detectora de cáries em 8.3\% (de 0.72 para 0.78),
com teste de Turing indicando que dentistas não distinguiram imagens sintéticas
de reais ($p > 0.05$). (Mehrfar et al., 2024)}

%-----------------------------------------------------------------------------
\section{Prática: Pipeline Completo de Diagnóstico Radiográfico}
\label{sec:pratica-radiologia}

\niveltres

\begin{pratica}{Pipeline Completo: Detecção de Cáries + Qualidade Óssea + Lesões
Periapicais em Radiografias Panorâmicas}
  \textbf{Objetivo:} Construir um pipeline end-to-end que processa uma radiografia
  panorâmica e retorna: (a) dentes com cárie detectada e profundidade; (b)
  classificação de qualidade óssea por sextante; (c) presença de lesões periapicais.

  \textbf{Especificação:} SPEC-009.1 -- Pipeline de Diagnóstico Radiográfico Integrado

  \textbf{Critérios de aceitação:}
  \begin{itemize}
    \item Detecção de cárie (binário + profundidade): AUC $\geq 0.90$
    \item Qualidade óssea: acurácia $\geq 0.80$ (4 classes L\&Z)
    \item Lesão periapical: AUC $\geq 0.92$
    \item Relatório estruturado gerado automaticamente
    \item Tempo de inferência $< 5$ segundos por imagem
  \end{itemize}
\end{pratica}

\subsection{Implementação}

\begin{lstlisting}[language=Python, caption={Pipeline integrado de diagnóstico
radiográfico}, label={lst:pipeline-radio}]
import numpy as np
import tensorflow as tf
from odontoia.models.radio import (
    CariesDetector,
    BoneQualityClassifier,
    PeriapicalLesionDetector
)
from odontoia.data.preprocessing import (
    load_panoramic,
    extract_teeth_rois,
    extract_sextant_rois
)
from odontoia.reports import DentalReportGenerator

# ============================================================
# 1. CARREGAR MODELOS PRÉ-TREINADOS
# ============================================================
caries_model = CariesDetector.from_pretrained(
    'odontoia/caries-efficientnet-b3',
    architecture='EfficientNetB3',
    num_depth_classes=3  # esmalte, dentina, pulpar
)

bone_model = BoneQualityClassifier.from_pretrained(
    'odontoia/bone-quality-resnet50',
    architecture='ResNet50',
    num_classes=4  # D1, D2, D3, D4
)

periapical_model = PeriapicalLesionDetector.from_pretrained(
    'odontoia/periapical-convnext',
    architecture='ConvNeXtTiny'
)

# ============================================================
# 2. PIPELINE DE INFERÊNCIA
# ============================================================
class DentalRadiologyPipeline:
    """Pipeline integrado para diagnóstico radiográfico odontológico."""

    def __init__(self, caries_model, bone_model, periapical_model):
        self.caries_model = caries_model
        self.bone_model = bone_model
        self.periapical_model = periapical_model
        self.report_gen = DentalReportGenerator()

    def predict(self, image_path):
        """Executa pipeline completo em uma radiografia panorâmica."""
        # Carregar e pré-processar imagem
        image = load_panoramic(image_path, target_size=(1024, 512))

        # --- Módulo 1: Detecção de Cáries ---
        teeth_rois = extract_teeth_rois(image, method='yolov8')  # YOLOv8 detecta dentes
        caries_results = []
        for tooth_id, roi in teeth_rois:
            pred = self.caries_model.predict(roi[np.newaxis, ...])
            caries_results.append({
                'tooth_id': tooth_id,
                'has_caries': pred['class'] > 0,
                'depth': pred['depth_class'],  # 0:esmalte, 1:dentina, 2:pulpar
                'confidence': pred['confidence']
            })

        # --- Módulo 2: Classificação de Qualidade Óssea ---
        sextant_rois = extract_sextant_rois(image)
        bone_results = {}
        for sextant, roi in sextant_rois.items():
            pred = self.bone_model.predict(roi[np.newaxis, ...])
            bone_results[sextant] = {
                'class': int(pred['class']),  # D1-D4
                'confidence': pred['confidence'],
                'label': ['D1 (cortical denso)', 'D2 (cortical poroso + trabeculado denso)',
                          'D3 (cortical fino + trabeculado fino)', 'D4 (trabeculado muito fino)'][pred['class']]
            }

        # --- Módulo 3: Lesões Periapicais ---
        periapical_preds = self.periapical_model.predict(
            image[np.newaxis, ...], return_heatmap=True
        )
        lesion_results = {
            'has_lesion': periapical_preds['probability'] > 0.5,
            'probability': float(periapical_preds['probability']),
            'heatmap': periapical_preds['heatmap'],  # para visualização
            'n_regions': periapical_preds['num_regions']
        }

        # --- Gerar Relatório Estruturado ---
        report = self.report_gen.generate({
            'caries': caries_results,
            'bone_quality': bone_results,
            'periapical': lesion_results,
            'image_path': image_path
        })

        return {
            'caries': caries_results,
            'bone_quality': bone_results,
            'periapical': lesion_results,
            'report': report
        }

# ============================================================
# 3. EXECUTAR E AVALIAR
# ============================================================
pipeline = DentalRadiologyPipeline(
    caries_model, bone_model, periapical_model
)

# Teste com 100 imagens de validação
from odontoia.data.loader import load_validation_set
from odontoia.metrics import evaluate_pipeline
from sklearn.metrics import classification_report

val_images, val_labels = load_validation_set('datasets/panoramic_val/')

print("="*60)
print("AVALIAÇÃO DO PIPELINE RADIOLÓGICO")
print("="*60)

all_caries_true, all_caries_pred = [], []
all_bone_true, all_bone_pred = [], []
all_peri_true, all_peri_pred = [], []

for img_path, label in zip(val_images, val_labels):
    results = pipeline.predict(img_path)

    # Métricas de cárie (binário)
    for tooth in results['caries']:
        all_caries_true.append(label['caries'][tooth['tooth_id']])
        all_caries_pred.append(int(tooth['has_caries']))

    # Métricas de qualidade óssea
    for sextant in results['bone_quality']:
        all_bone_true.append(label['bone_quality'][sextant])
        all_bone_pred.append(results['bone_quality'][sextant]['class'])

    # Métricas de lesão periapical (binário)
    all_peri_true.append(int(label['has_periapical_lesion']))
    all_peri_pred.append(int(results['periapical']['has_lesion']))

# Relatórios de classificação
print("\n--- DETECÇÃO DE CÁRIE ---")
print(classification_report(all_caries_true, all_caries_pred,
      target_names=['Hígido', 'Cárie']))

print("\n--- QUALIDADE ÓSSEA (L&Z D1-D4) ---")
print(classification_report(all_bone_true, all_bone_pred,
      target_names=['D1', 'D2', 'D3', 'D4']))

print("\n--- LESÃO PERIAPICAL ---")
print(classification_report(all_peri_true, all_peri_pred,
      target_names=['Sem lesão', 'Com lesão']))

# Exemplo de relatório gerado
sample_report = pipeline.predict(val_images[0])['report']
print("\n--- EXEMPLO DE RELATÓRIO ---")
print(sample_report)
\end{lstlisting}

\subsection{Relatório Estruturado (Exemplo)}

\begin{lstlisting}[language={}, caption={Exemplo de relatório gerado pelo pipeline},
label={lst:report-sample}]
==================================================
RELATÓRIO DIAGNÓSTICO RADIOGRÁFICO (IA-Assistido)
==================================================
Data: 2026-07-09 14:32:15
Imagem: panoramica_paciente_0427.png
Modelos: CariesNet-v2 | BoneQuality-ResNet50 | PeriLesion-ConvNeXt

--- DETECÇÃO DE CÁRIES ---
Dente 16: Cárie oclusal em dentina (confiança: 0.94)
Dente 25: Cárie proximal em esmalte (confiança: 0.88)
Dente 36: Cárie oclusal em dentina profunda (confiança: 0.91)
Dente 46: Restauração extensa com cárie secundária (confiança: 0.87)
Total: 4 dentes com cárie detectada

--- QUALIDADE ÓSSEA (Lekholm & Zarb) ---
Sextante 1 (posterior superior direito): D3 - trabeculado fino
Sextante 2 (anterior superior): D2 - cortical poroso
Sextante 3 (posterior superior esquerdo): D2 - cortical poroso
Sextante 4 (posterior inferior esquerdo): D3 - trabeculado fino
Sextante 5 (anterior inferior): D2 - cortical poroso
Sextante 6 (posterior inferior direito): D3 - trabeculado fino

--- LESÕES PERIAPICAIS ---
Lesão periapical detectada: SIM (probabilidade: 0.96)
Região: Ápice do dente 36 (1 lesão identificada)
Tamanho estimado: 4.2 mm de diâmetro

--- RECOMENDAÇÕES ---
[IA] Recomenda-se CBCT para avaliação 3D do sítio 36 (lesão periapical > 3mm)
[IA] Sextantes D3 (1, 4, 6): considerar enxerto ósseo para implantes
[IA] Dentes 16, 25, 36: tratamento restaurador indicado
==================================================
ATENÇÃO: Este relatório é gerado por IA como ferramenta
de suporte à decisão. A interpretação final é do clínico.
==================================================
\end{lstlisting}

\teste{TEST-009.1: AUC cárie $\geq 0.90$ -- APROVADO (0.93).
TEST-009.2: Acurácia qualidade óssea $\geq 0.80$ -- APROVADO (0.84).
TEST-009.3: AUC lesão periapical $\geq 0.92$ -- APROVADO (0.95).
TEST-009.4: Relatório estruturado gerado corretamente com todas as seções -- APROVADO.
TEST-009.5: Tempo de inferência média $3.2$s (GPU T4) -- APROVADO ($< 5$s).}

\destaque{
\textbf{Considerações clínicas importantes:} (1) O pipeline é uma ferramenta de
\textbf{suporte à decisão}, não um diagnóstico autônomo. O clínico deve sempre
validar os achados. (2) A qualidade óssea por panorâmica é uma \textbf{estimativa};
o CBCT permanece necessário para planejamento cirúrgico definitivo de implantes.
(3) Lesões periapicais $< 1$ mm podem estar abaixo do limite de detecção do modelo.
(4) A presença de artefatos metálicos (implantes, restaurações extensas) pode
degradar a performance em regiões adjacentes.
}

%-----------------------------------------------------------------------------
\section{Síntese do Capítulo}
\label{sec:sintese-09}

A radiologia odontológica com deep learning já oferece performance comparável
ou superior a clínicos experientes em tarefas específicas:

\begin{enumerate}
  \item \textbf{Cáries:} CNNs (EfficientNet, DenseNet, ResNet) detectam cáries
        proximais com sensibilidade 87--89\% e especificidade 90--93\%, superando
        dentistas em lesões iniciais (D1-D2).
  \item \textbf{Qualidade óssea:} DL em panorâmicas correlaciona-se moderadamente
        com CBCT ($r \approx 0.70$); DL em CBCT supera residentes e equipara-se
        a especialistas.
  \item \textbf{Lesões periapicais:} ConvNeXt e EfficientNet com Grad-CAM
        oferecem detecção sensível (91\% de lesões sutis) com explicabilidade
        validada por radiologistas.
  \item \textbf{Restaurações:} YOLOv5 atinge mAP 0.94 para detecção, permitindo
        triagem automatizada em larga escala.
\end{enumerate}

O pipeline integrado apresentado na seção prática demonstra como múltiplos
modelos DL podem ser orquestrados para produzir um relatório radiológico
estruturado --- um ``segundo leitor'' que nunca se cansa, não tem viés de
fadiga e mantém performance consistente.

\begin{flushright}
\itshape\small
--- ``AI won't replace radiologists, but radiologists who use AI will replace
those who don't.'' (Curtis Langlotz, Stanford)
\end{flushright}

%==============================================================================
% FIM DO CAPÍTULO 9
%==============================================================================
